\section{\ENG{Git}, \ENG{GitHub} και Πρακτική Ελέγχου Εκδόσεων}
\label{sec:appendix_versioning_governance}

Το οικοσύστημα δεν υποστηρίζεται μόνο από διάσπαση σε διαφορετικά \ENG{modules}, αλλά και από μια σαφή,
αν και όχι πλήρως ομοιόμορφη, πρακτική ελέγχου εκδόσεων. Τα \ENG{tags} των αποθετηρίων τεκμηριώνουν πραγματική λογική έκδοσης
και όχι απλή αποθήκευση πρόχειρου κώδικα.
Τη στιγμή που γράφονται αυτές οι γραμμές, το \texttt{\ENG{scorm-engine}}, το \texttt{\ENG{scorm-engine-php-sdk}}
και το \texttt{\ENG{example-lms-client}} φέρουν \ENG{tag} \texttt{\ENG{v1.0.0}}, ο \texttt{\ENG{player}} έχει ήδη προχωρήσει
έως το \texttt{\ENG{v1.0.1}} και το \texttt{\ENG{central-docker-infrastructure}} έως το \texttt{\ENG{v1.0.2}}.
Η ύπαρξη αυτών των \ENG{snapshots} συνδέεται ευθέως με αυτοματοποιημένες ροές δημιουργίας έκδοσης (κυκλοφορίας) και δημιουργίας \ENG{changelogs},
οι οποίες αναλύονται λεπτομερέστερα στο επόμενο παράρτημα.

Η \ENG{GitHub}-side διαχείριση είναι ελαφρά, αλλά παρούσα. Τα αρχεία \texttt{\ENG{CODEOWNERS}} στα περισσότερα βασικά αποθετήρια δηλώνουν
βασική ιδιοκτησία του κώδικα περιορίζοντας τη δυνατότητα αλλαγών από μη εξουσιοδοτημένους χρήστες,
ενώ τα \texttt{\ENG{labels.json}}, \texttt{\ENG{labeler-config.yml}} και τα αντίστοιχα \ENG{workflows} υπάρχουν για να βοηθήσουν 
στην αυτοματοποίηση ορισμένων διαδικασιών. Η πρακτική αυτή δεν συνιστά πλήρες \ENG{repository governance framework},
αλλά είναι αρκετή ώστε να τεκμηριώσει ότι ο οργανισμός δεν λειτουργεί ως άναρχη συλλογή φακέλων.

Για λόγους ομοιομορίας και αυτοματοποίησης κοινών διαδικασιών, το \texttt{\ENG{github-actions-template}} 
περιέχει τοπικό \ENG{wizard} που μπορεί να δημιουργήσει αυτά τα \ENG{GitHub Actions} με ομοιόμορφο τρόπο.


\begin{table}[H]
\centering
\caption{Ενδείξεις ελέγχου εκδόσεων, \ENG{release} πρακτικής και ελαφράς διακυβέρνησης στα βασικά αποθετήρια}
\label{tab:appendix_versioning_workflow_conventions}
\begingroup
\footnotesize
\setlength{\tabcolsep}{3pt}
\renewcommand{\arraystretch}{1.20}
\begin{adjustbox}{center,max width=\textwidth}
\begin{tabular}{|p{2.8cm}|p{2.2cm}|p{2.2cm}|p{4.0cm}|p{4.1cm}|}
\hline
Αποθετήριο & Κύριος κλάδος που τεκμηριώνεται & \ENG{Tags} που παρατηρούνται & \ENG{Governance} και \ENG{workflow artifacts} & Παρατήρηση \\
\hline
\texttt{\ENG{scorm-engine}} & \texttt{\ENG{master}} & \texttt{\ENG{v1.0.0}} & \texttt{\ENG{CODEOWNERS}}, \ENG{labeler}, \ENG{semantic-release}, \ENG{changelog workflow} & Η \ENG{release} λογική και οι \ENG{test workflows} είναι συνεπείς με τον κύριο κλάδο \\
\hline
\texttt{\ENG{scorm-engine-php-sdk}} & \texttt{\ENG{master}} τοπικά, \texttt{\ENG{main/master}} σε workflows & \texttt{\ENG{v1.0.0}} & \texttt{\ENG{CODEOWNERS}}, \ENG{labeler}, \ENG{release workflow}, \ENG{tests workflow} & Υπάρχει \ENG{branch} ασυνέπεια που απαιτεί προσεκτική διατύπωση \\
\hline
\texttt{\ENG{example-lms-client}} & \texttt{\ENG{master}} & \texttt{\ENG{v1.0.0}} & \texttt{\ENG{CODEOWNERS}}, \ENG{labeler}, \ENG{release} και \ENG{changelog workflows} & Η \ENG{release} ιστορία συνδέεται με αυτοματοποιημένη ενημέρωση \ENG{CHANGELOG} \\
\hline
\texttt{\ENG{player}} & \texttt{\ENG{master}} & \texttt{\ENG{v1.0.0}}, \texttt{\ENG{v1.0.1}} & \ENG{labeler}, \ENG{semantic-release}, \ENG{tests/lint workflows} & Δεν τεκμηριώνεται \texttt{\ENG{CODEOWNERS}} αρχείο, αλλά υπάρχει πλήρες \ENG{workflow set} \\
\hline
\texttt{\ENG{central-docker-infrastructure}} & \texttt{\ENG{master}} & \texttt{\ENG{v1.0.0}} έως \texttt{\ENG{v1.0.2}} & Προσαρμοσμένο \ENG{release workflow}, \ENG{labeler}, \ENG{CI} και \ENG{lint workflows} & Η \ENG{release} πρακτική είναι χειροποίητη και όχι \ENG{semantic-release}-based \\
\hline
\end{tabular}
\end{adjustbox}
\endgroup

\end{table}
